\chapter{Present the Work}\label{ch:talks}

A talk is an argument the audience hears once, in order, from the person who presents it.
The agent can draft the deck, but the presenter must be able to defend every frame, so
the argument is settled with them before any slide is drawn.

Beamer, the \LaTeX\ class used here, calls a slide a frame. A frame that reveals its
content step by step with \code{\\pause} becomes several PDF pages, one per overlay,
which share one page label. The render gate and the review inbox are those of
Chapter~\ref{ch:reports}.

\section{Settle the storyline first}

\begin{headline}
  The agent writes the storyline and the presenter confirms it. Slides are built only
  from a confirmed storyline, and later changes to the argument go back to the presenter.
\end{headline}

\begin{table}[H]
  \small
  \begin{tabularx}{\textwidth}{L{0.2\textwidth}Y}
    \toprule
    Field & Content \\
    \midrule
    Kind & Own work, re-cut, someone else's paper, group-meeting update or poster. \\
    \rowcolor{ctpSurface0!30}
    Audience & Who listens, what they know, what they will do with the talk. \\
    Slot & Length of the talk and of the questions. \\
    \rowcolor{ctpSurface0!30}
    Question & The one question the talk answers. \\
    Take-away & One sentence the audience should repeat afterwards. \\
    \rowcolor{ctpSurface0!30}
    Slides & One line per frame: its title as a claim, and a time mark. \\
    Backup & Expected questions, each with the backup frame that answers it. \\
    \rowcolor{ctpSurface0!30}
    Sources & Where every figure, number and reused diagram comes from. \\
    Formats & Audience PDF and presenter copy; a \code{.pptx} if the room needs PowerPoint. \\
    \bottomrule
  \end{tabularx}
  \caption{The storyline contract \cite[The storyline contract]{skill:deck-production}.
    Choices the agent made alone are listed in it as agent defaults.}
  \label{tab:talks-storyline}
\end{table}

\begin{figure}[H]
  \centering
  \begin{tikzpicture}[node distance=6mm and 8mm]
    \node[work,text width=30mm] (read) {Read the report, plotbook, paper or old deck};
    \node[work,text width=30mm,right=of read] (story) {Write the storyline};
    \node[ask,right=of story] (ok) {Confirmed?};
    \node[work,text width=30mm,below=10mm of ok] (build) {Build deck and notes from the
      template};
    \node[open,text width=30mm,left=of build] (teach) {Presenter walks through the frames};
    \node[check,text width=30mm,left=of teach] (gate) {Slides gate, visual and copy
      review};
    \node[result,text width=30mm,below=16mm of gate] (inbox) {PDFs and \code{.pptx} in the
      review inbox};
    \draw[flow] (read)--(story);
    \draw[flow] (story)--(ok);
    \draw[back] (ok.north) to[out=135,in=45] node[lbl,above]{no: revise} (story.north);
    \draw[flow] (ok)--node[lbl,right]{yes}(build);
    \draw[flow] (build)--(teach);
    \draw[flow] (teach)--(gate);
    \draw[back] (gate.south east) to[out=-25,in=-155] node[lbl,below]{fix in the source}
    (build.south west);
    \draw[flow] (gate)--(inbox);
  \end{tikzpicture}
  \caption{From source material to a deck awaiting review \cite[Storyline first; Build;
      Check and deliver]{skill:presentations}.}
  \label{fig:talks-flow}
\end{figure}

\section{Build the deck}

\begin{headline}
  Each frame title states a claim, the figures are the reviewed plotbook files, the diagrams
  are the report's own asset files, and the talk adds no claim the written work does not
  support.
\end{headline}

\begin{table}[H]
  \small
  \begin{tabularx}{\textwidth}{L{0.24\textwidth}L{0.18\textwidth}Y}
    \toprule
    Kind of talk & Size & What changes \\
    \midrule
    Own work (defence, colloquium, seminar) & set by the slot & Depth follows the audience;
    derivations and software details move to backup frames. \\
    \rowcolor{ctpSurface0!30}
    Re-cut of a deck & set by the slot & A new storyline; frames kept, moved to backup or removed are
    recorded. \\
    Someone else's paper & set by the slot & Claims kept apart from your assessment; a
    strengths and limits section; a rubric map when graded. \\
    \rowcolor{ctpSurface0!30}
    Group-meeting update & 3--8 frames & What changed, the plots that show it, the decisions
    needed; review status on the slides. \\
    Poster & A0 portrait & Blocks for question, method, result and take-away; larger font floors. \\
    \bottomrule
  \end{tabularx}
  \caption{Kinds of talk \cite[Kinds of talk]{skill:deck-production}.}
  \label{tab:talks-kinds}
\end{table}

Slide rules \cite[Slide rules]{skill:deck-production}:
\begin{itemize}
  \item one idea per frame; the sentences go in the speaker notes;
  \item \code{\\pause} only where the audience must see the problem before the answer;
  \item a \code{\\scriptsize} source line under every borrowed figure or number;
  \item backup frames linked from the frame that raises the question, and back;
  \item the report's diagrams input from the study's \code{assets/} files with the same number
        macros, and the plotbook's position strip in the title bar of frames about one step
        \cite[Diagrams and navigation]{skill:deck-production}.
\end{itemize}

The template is a 16:9 Beamer deck and an A0 poster with the palette of the reports. One
source builds the audience PDF and the presenter copy, whose right half holds the speaker
notes.

\section{Check and deliver}

\begin{headline}
  The presenter walks through the deck, then it passes the slides gate, visual review and,
  for a judged talk, copy review before it reaches the review inbox.
\end{headline}

\begin{table}[H]
  \small
  \begin{tabularx}{\textwidth}{L{0.3\textwidth}Y}
    \toprule
    Slides profile & Why \\
    \midrule
    Overlays grouped into frames by page label & Text hidden by \code{\\pause} is not a
    finding until the frame's last overlay. \\
    \rowcolor{ctpSurface0!30}
    No whitespace or density checks & An empty slide area is a design choice. \\
    Small text inside a placed figure is flagged & Re-rendering a plotbook figure is the
    deck owner's choice; text set on the slide keeps the fail floor. \\
    \rowcolor{ctpSurface0!30}
    Text past the slide edge fails & The audience cannot read it. \\
    \bottomrule
  \end{tabularx}
  \caption{What the slides profile changes \cite[Slides]{skill:render-workflow}.}
  \label{tab:talks-gate}
\end{table}

\begin{itemize}
  \item \review\ Before the gate the presenter explains back every frame they did not
        draft; gaps become notes or backup frames, and an absent presenter's walk-through
        is recorded as pending \cite[Notes and
          questions]{skill:deck-production}.
  \item \evidence\ The PowerPoint copy, when the room needs one, holds one picture per
        slide and the speaker notes; edit the source, not the \code{.pptx}.
  \item \gap\ Uploading to a meeting page or sending to the audience is the presenter's
        step, never the agent's.
\end{itemize}

\hook{talks-delivery}
